\section{Introduction}
This chapter reports the evaluation of VeriGate: functional and test verification, the headline correctness result (false-hit elimination) with its ablation, the mechanism by failure mode, the Tier-2 NLI trade-off, the latency and vector-index results, the backend comparison, the cost-aware cascade, the security (tenant-isolation) result, the live LLM demonstration, observability, and the operator dashboard. Every figure and table is regenerable from the repository, and the evaluation harness shares the exact \texttt{CachePolicy} used by the live gateway, so measured behaviour equals shipped behaviour.

\subsection{Test-Suite Verification}
The automated suite is the first line of evidence that each component behaves as specified. All \textbf{18 tests pass} in the reference environment (Python~3.13, Windows~11).

\begin{table}[H]
    \centering
    \caption{Representative automated tests and the guarantee each proves.}
    \label{tab:tests}
    \footnotesize
    \begin{tabular}{|l|p{8.2cm}|l|}
    \hline
    \textbf{Test} & \textbf{Guarantee proven} & \textbf{Result} \\ \hline
    test\_health & server boots and reports live configuration & PASS \\ \hline
    test\_auth\_required & requests without a valid key are rejected (401) & PASS \\ \hline
    test\_cache\_miss\_then\_hit & a paraphrase is served from cache after the first miss & PASS \\ \hline
    test\_verifier\_rejects\_number\_mismatch & a look-alike differing by a number is NOT served & PASS \\ \hline
    test\_volatile\_query\_bypasses\_cache & time-sensitive queries are never cached & PASS \\ \hline
    test\_rate\_limit\_returns\_429 & a burst is throttled & PASS \\ \hline
    test\_tenant\_isolation & tenant B cannot read tenant A's cached answer & PASS \\ \hline
    test\_cascade\_routing & easy$\rightarrow$cheap, hard$\rightarrow$strong routing is correct & PASS \\ \hline
    test\_budget\_blocks\_over\_limit & over-budget calls are 429; cache hits stay free & PASS \\ \hline
    test\_metrics\_endpoint & Prometheus metrics are exposed & PASS \\ \hline
    \end{tabular}
\end{table}

\section{Headline Result: Eliminating False Hits}
The primary success metric is the false-hit rate --- how often the cache serves a wrong answer to a look-alike query. VeriGate was compared against a fixed-threshold baseline (GPTCache-style, $T=0.72$) on the curated benchmark, adding one component at a time. Table~\ref{tab:ablation} gives the ablation.

\begin{table}[H]
\centering
\caption{Ablation on the benchmark. VeriGate serves zero wrong answers (100\% precision) versus the baseline's 46.7\% false-hit rate.}
\label{tab:ablation}
\begin{tabular}{|l|c|c|c|c|}
\hline
\textbf{System} & \textbf{Hit-rate} & \textbf{False-hit rate} & \textbf{Precision} & \textbf{F1} \\ \hline
Baseline (static $T=0.72$)       & 40.0\% & 46.7\% & 46.2\% & 0.43 \\ \hline
\ \ + Adaptive threshold         & 33.3\% & 33.3\% & 50.0\% & 0.40 \\ \hline
\ \ + Verifier (Tier-1)          & 40.0\% & 13.3\% & 75.0\% & 0.52 \\ \hline
\textbf{VeriGate (full)}         & 33.3\% & \textbf{0.0\%} & \textbf{100.0\%} & 0.50 \\ \hline
\end{tabular}
\end{table}

The baseline serves a wrong answer on nearly half of the hard negatives. VeriGate serves \textbf{none}: it is stricter, so its raw hit-rate is slightly lower, which is the intended and correct trade-off for a system whose priority is correctness. Figure~\ref{fig:tradeoff} makes the trade-off visual.

\begin{figure}[H]
    \centering
    \includegraphics[width=0.82\textwidth,height=0.32\textheight,keepaspectratio]{tradeoff.png}
    \caption{The trade-off (``money'') curve. Sweeping the similarity threshold, false-hit rate (x, lower is better) is plotted against hit-rate (y, higher is better) for the baseline (red) and the verifier-enabled system (blue); the green star is full VeriGate. The blue curve lies \emph{above and to the left} of the red one --- at any hit-rate it serves far fewer wrong answers (Pareto domination).}
    \label{fig:tradeoff}
\end{figure}

\begin{figure}[H]
    \centering
    \includegraphics[width=0.82\textwidth,height=0.30\textheight,keepaspectratio]{ablation.png}
    \caption{Ablation bars: hit-rate (green) and false-hit rate (red) as components are added. False-hit rate falls step-by-step to zero for full VeriGate while hit-rate stays reasonable.}
    \label{fig:ablation}
\end{figure}

\subsection{False Hits by Failure Mode (the Mechanism)}
Table~\ref{tab:modes} decomposes the false hits by trap type, showing precisely which component neutralises each. The \emph{verifier} eliminates entity-, number-, and negation-swap false hits; the \emph{adaptive threshold} clears the related-topic ones; together they reach zero.

\begin{table}[H]
    \centering
    \caption{False hits by failure mode (lower is better). Cells show false hits / traps of that type.}
    \label{tab:modes}
    \begin{tabular}{|l|c|c|c|c|}
    \hline
    \textbf{System} & \textbf{entity\_swap} & \textbf{number\_swap} & \textbf{negation} & \textbf{related\_topic} \\ \hline
    Baseline           & 0/5 & 3/3 & 2/2 & 2/5 \\ \hline
    \ \ + Adaptive     & 0/5 & 3/3 & 2/2 & 0/5 \\ \hline
    \ \ + Verifier     & 0/5 & 0/3 & 0/2 & 2/5 \\ \hline
    \textbf{VeriGate}  & 0/5 & \textbf{0/3} & \textbf{0/2} & \textbf{0/5} \\ \hline
    \end{tabular}
\end{table}

\subsection{Volatility (Staleness) Detector}
On the volatility set the detector achieved \textbf{5/5 recall} on volatile queries (all correctly bypassed) with \textbf{0/4 false positives} on stable controls, removing the staleness class of errors entirely.

\section{Tier-2 NLI: A Correctness--Recall Trade-off}
Adding the optional Tier-2 NLI verifier was evaluated in isolation (Table~\ref{tab:nli}).

\begin{table}[H]
    \centering
    \caption{Effect of adding Tier-2 NLI to the full system.}
    \label{tab:nli}
    \begin{tabular}{|l|c|c|c|c|}
    \hline
    \textbf{System} & \textbf{Hit-rate} & \textbf{False-hit} & \textbf{Precision} & \textbf{F1} \\ \hline
    VeriGate (Tier-1)           & 33.3\% & 0.0\% & 100.0\% & 0.50 \\ \hline
    VeriGate (Tier-1 + Tier-2)  & 20.0\% & 0.0\% & 100.0\% & 0.33 \\ \hline
    \end{tabular}
\end{table}

Tier-2 achieves zero false hits at \emph{every} threshold but caps recall lower, because NLI is conservative and rejects some genuine paraphrases. Crucially, it is \textbf{largely redundant once the adaptive threshold is active} --- both target the same related-topic negatives. The design conclusion, a defensible engineering decision rather than a naive ``add more components'' story, is: \emph{the recommended default is adaptive + Tier-1 (zero false hits at higher recall), with Tier-2 offered as an opt-in ``max-correctness mode''} for correctness-critical domains (medical, legal, finance) or when the base threshold must be kept loose.

\section{Latency and the Cache-Hit Path}
Per-stage latency was measured to answer the direct question: \emph{how much does the verification that constitutes the contribution actually cost?}

\begin{table}[H]
    \centering
    \caption{Per-stage latency (milliseconds). Tier-1 verification is effectively free; Tier-2 NLI is a deliberate opt-in cost.}
    \label{tab:latency}
    \begin{tabular}{|l|c|c|c|}
    \hline
    \textbf{Stage} & \textbf{p50} & \textbf{p95} & \textbf{p99} \\ \hline
    Embed query                    & 19.0  & 30.3  & 84.4 \\ \hline
    KNN lookup --- 100 entries     & 27.1  & 35.9  & 41.2 \\ \hline
    KNN lookup --- 1000 entries    & 292.9 & 346.8 & 422.8 \\ \hline
    \textbf{Tier-1 verify}         & \textbf{0.03} & 0.03 & 0.05 \\ \hline
    \textbf{Tier-2 NLI verify}     & \textbf{171.7} & 183.5 & 187.9 \\ \hline
    HIT path (embed+KNN@1000+Tier-1) & 334.3 & 374.0 & 399.3 \\ \hline
    \end{tabular}
\end{table}

\textbf{RQ (verification overhead):} Tier-1 costs $\sim$0.03~ms --- structural safety is effectively free on the hit path. Tier-2 NLI adds $\sim$172~ms, a second reason it is opt-in. Table~\ref{tab:latency} also exposes the real bottleneck in the naive design: the brute-force $O(n)$ KNN scan, which grows to $\sim$293~ms at 1000 entries.

\begin{figure}[H]
    \centering
    \includegraphics[width=0.82\textwidth,height=0.30\textheight,keepaspectratio]{latency.png}
    \caption{Per-stage latency (log scale) against an 800~ms hosted-LLM reference. Tier-1 verification is invisible ($\sim$free); the costly stages are the NLI check and the KNN scan, both still below the LLM reference.}
    \label{fig:latency}
\end{figure}

\subsection{Vector-Index Upgrade}
The $O(n)$ scan was replaced by a vectorised in-process index (one NumPy matrix, a single matrix--vector product). Table~\ref{tab:index} shows the effect.

\begin{table}[H]
    \centering
    \caption{KNN lookup latency before and after the vectorised index.}
    \label{tab:index}
    \begin{tabular}{|l|c|c|c|}
    \hline
    \textbf{Cache size} & \textbf{Before (parse loop)} & \textbf{After (vectorised)} & \textbf{Speedup} \\ \hline
    1000 entries & $\sim$293 ms & \textbf{0.04 ms} & $\sim$7{,}000$\times$ \\ \hline
    5000 entries & $\sim$1.5 s (extrap.) & \textbf{0.11 ms} & --- \\ \hline
    \end{tabular}
\end{table}

After the upgrade the full cache-hit path is \textbf{$\sim$10~ms} --- about \textbf{78$\times$} faster than an 800~ms LLM call --- and is now bounded by the embedding step, not the lookup. Figure~\ref{fig:scaling} shows lookup staying sub-millisecond and near-flat as the cache grows.

\begin{figure}[H]
    \centering
    \includegraphics[width=0.82\textwidth,height=0.30\textheight,keepaspectratio]{latency_scaling.png}
    \caption{KNN lookup latency versus cache size. The vectorised in-process index stays sub-millisecond and near-flat to 5000 entries, versus $\sim$293~ms at 1000 for the original parse loop.}
    \label{fig:scaling}
\end{figure}

\subsection{Backend Comparison (In-Process vs RediSearch)}
The RediSearch/HNSW backend was verified end-to-end on a live Redis Stack (a semantic paraphrase HIT and the Austria/Australia trap rejected, both served from the shared index). Table~\ref{tab:backend} compares raw latency.

\begin{table}[H]
    \centering
    \caption{KNN latency (p50): in-process NumPy versus RediSearch/HNSW.}
    \label{tab:backend}
    \begin{tabular}{|l|c|c|}
    \hline
    \textbf{Cache size} & \textbf{In-process (NumPy)} & \textbf{RediSearch (HNSW)} \\ \hline
    1000  & 0.06 ms & 3.60 ms \\ \hline
    5000  & 0.15 ms & 3.83 ms \\ \hline
    10000 & 0.24 ms & 4.69 ms \\ \hline
    \end{tabular}
\end{table}

\begin{figure}[H]
    \centering
    \includegraphics[width=0.80\textwidth,height=0.30\textheight,keepaspectratio]{backend_compare.png}
    \caption{In-process versus RediSearch KNN latency. The in-process index wins on raw latency (no network hop); RediSearch carries a $\sim$3.6~ms round-trip floor but stays nearly flat as the cache grows and is shared across replicas. Both are two-to-four orders of magnitude faster than an LLM call. Recommendation: in-process for a single node, RediSearch for horizontal scale.}
    \label{fig:backend}
\end{figure}

\section{Cost-Aware Cascade}
The second contribution was evaluated by measuring cost reduction against quality retention (cheap cost = 1, strong cost = 15$\times$). Table~\ref{tab:cascade} and Figure~\ref{fig:cascade} give the result.

\begin{table}[H]
    \centering
    \caption{Cascade cost versus quality across routing thresholds.}
    \label{tab:cascade}
    \begin{tabular}{|c|c|c|c|c|c|}
    \hline
    \textbf{Threshold} & \textbf{Escalated} & \textbf{Cost reduction} & \textbf{Quality retained} & \textbf{vs random} & \textbf{Routing F1} \\ \hline
    0.30            & 50\% & \textbf{47\%} & \textbf{100\%} & 50\% & \textbf{1.00} \\ \hline
    0.35 (default)  & $\sim$45\% & $\sim$52\% & $\sim$95\% & $\sim$45\% & $\sim$0.95 \\ \hline
    0.50            & 42\% & 54\% & 85\% & 42\% & 0.92 \\ \hline
    0.70            & 18\% & 77\% & 35\% & 18\% & 0.52 \\ \hline
    \end{tabular}
\end{table}

At threshold~0.30 the cascade cuts LLM cost by \textbf{$\sim$47\%} with \textbf{100\% quality retained} on hard queries, and it sits above random routing at every cost level --- a hard query is far likelier to be escalated than under random routing.

\begin{figure}[H]
    \centering
    \includegraphics[width=0.80\textwidth,height=0.30\textheight,keepaspectratio]{cascade.png}
    \caption{Quality retention versus cost reduction. The complexity classifier holds 100\% quality up to $\sim$47\% cost reduction, then declines, and dominates random routing at every operating point.}
    \label{fig:cascade}
\end{figure}

\section{Security: Cross-Tenant Isolation}
The threat model identified cross-tenant cache leakage as the single most serious risk of a semantic cache. After partitioning the cache per tenant, tenant~B asking tenant~A's \emph{exact} question receives a \textbf{MISS} (its own empty cache), while A re-asking receives a HIT --- verified in-process by \texttt{test\_tenant\_isolation} and, on the RediSearch backend, by a smoke-test assertion that \texttt{t\_other} cannot retrieve \texttt{t\_demo}'s entry. The defining vulnerability of the design was thus identified and then eliminated with a regression test guarding it.

\section{Live LLM Demonstration}
With a real Gemini key, the gateway returned genuine Gemini answers and then served repeats and paraphrases from cache (Table~\ref{tab:live}). This is the whole thesis on a real model: a slow, paid call becomes an instant, free cache hit --- even for a reworded question.

\begin{table}[H]
    \centering
    \caption{Live gateway with Gemini: first call versus subsequent cache hits.}
    \label{tab:live}
    \begin{tabular}{|l|l|c|p{5.2cm}|}
    \hline
    \textbf{Cache} & \textbf{Provider} & \textbf{Latency} & \textbf{Query} \\ \hline
    MISS & gemini & 23{,}048 ms & ``what is a semantic cache'' (real Gemini call) \\ \hline
    HIT  & cache  & 22 ms & ``what is a semantic cache'' ($\sim$1000$\times$ faster, \$0) \\ \hline
    HIT  & cache  & 20 ms & ``explain what a semantic cache is'' (paraphrase matched by meaning) \\ \hline
    \end{tabular}
\end{table}

\section{Observability}
The provisioned Prometheus + Grafana stack was verified live: the Prometheus target was healthy and Grafana rendered the cache-hit-rate gauge at \textbf{51.1\%} (24 hits / 47 requests) alongside request-rate, cache-outcome, and p95-latency panels driven by real traffic --- demonstrating that the same metrics used for evaluation are exposed for production monitoring.

\section{Operator Dashboard (User Interface) Results}
The gateway serves a single-page operator dashboard at \texttt{/ui} for interactive testing. Figure~\ref{fig:ui_hero} shows the landing view; Figure~\ref{fig:ui_hit} shows a paraphrase served as a HIT; and Figure~\ref{fig:ui_blocked} shows the single clearest proof of the contribution --- a look-alike whose similarity (0.944) is \emph{above} the threshold (0.791) yet is rejected with \texttt{verify\_rejected:number\_mismatch}.

\begin{figure}[H]
    \centering
    \includegraphics[width=0.90\textwidth]{ui_hero.png}
    \caption{VeriGate operator dashboard --- landing view. The header reports the live configuration (in-process index, mock provider).}
    \label{fig:ui_hero}
\end{figure}

\begin{figure}[H]
    \centering
    \includegraphics[width=0.90\textwidth]{ui_hit.png}
    \caption{A reworded query (``what is the process to recover my account password'') served as a HIT from a prior ``how do I reset my password'' --- matched by meaning, in $\sim$30~ms.}
    \label{fig:ui_hit}
\end{figure}

\begin{figure}[H]
    \centering
    \includegraphics[width=0.90\textwidth]{ui_blocked.png}
    \caption{The self-verifier in action. The ``50 gb'' query is similar (0.944) to a cached ``5 gb'' query --- \emph{above} the adaptive threshold (0.791) --- yet is BLOCKED with reason \texttt{verify\_rejected:number\_mismatch}. This is precisely the false hit a fixed-threshold cache would serve, and precisely what VeriGate prevents.}
    \label{fig:ui_blocked}
\end{figure}

\section{Summary of Results}
VeriGate reduces the false-hit rate from 46.7\% to 0.0\% (100\% precision) at a structural-verification cost of $\sim$0.03~ms; a vectorised index brings the cache-hit path to $\sim$10~ms ($\sim$78$\times$ faster than an LLM call) and keeps lookup sub-millisecond as the cache grows; the cost-aware cascade cuts LLM spend $\sim$47\% with no quality loss on hard queries; and the system is verified secure (tenant isolation), reliable (failover, breaker), observable (Grafana), and live (real Gemini). Collectively these results substantiate every objective of Section~1.4.
